% !PRINTED 69
\weberpage{69}
%<<CONT

soll \(\Delta\) sein. \(\Delta\) wird die \gesperrt{Determinante} der \(n^2\) Elemente \(a_i^{(k)}\) genannt, und zwar, wenn die Unterscheidung nothwendig ist, eine \gesperrt{\(n\)-reihige Determinante} (auch Determinante \(n\)ten Grades oder \(n\)ter Ordnung). Das Glied \(M\) dieser Summe, d.~h. also das Product aller in der Diagonale des Quadrats stehenden Elemente, wird das \gesperrt{Hauptglied} genannt.

Nehmen wir z.~B. \(n=2\), so erhalten wir
\begin{equation}
\Delta=a_1^{(1)}a_2^{(2)}-a_2^{(1)}a_1^{(2)},
\tag{4}
\end{equation}
und für \(n=3\) [nach (3) des vorigen Paragraphen]:
\begin{equation}
\begin{aligned}
\Delta={}&a_1^{(1)}a_2^{(2)}a_3^{(3)}
+a_2^{(1)}a_3^{(2)}a_1^{(3)}
+a_3^{(1)}a_1^{(2)}a_2^{(3)}\\
&-a_1^{(1)}a_3^{(2)}a_2^{(3)}
-a_2^{(1)}a_1^{(2)}a_3^{(3)}
-a_3^{(1)}a_2^{(2)}a_1^{(3)}.
\end{aligned}
\tag{5}
\end{equation}
Oder in anderer Bezeichnung:
\begin{equation}
\begin{vmatrix}
a&b\\
c&d
\end{vmatrix}
=ad-bc,
\tag{6}
\end{equation}
\begin{equation}
\begin{vmatrix}
a&b&c\\
a'&b'&c'\\
a''&b''&c''
\end{vmatrix}
=ab'c''+bc'a''+ca'b''-ac'b''-ba'c''-cb'a''.
\tag{7}
\end{equation}

Es ist dem Leser zu empfehlen, die Berechnung solcher Determinanten an Zahlenbeispielen einzuüben.

Die Bezeichnung (1) ist in vielen Fällen zu umständlich; es sind daher noch andere, kürzere Zeichen in Gebrauch. So setzt \gesperrt{Jacobi}, indem er nur das Hauptglied der entwickelten Determinante ausführlich schreibt:
\begin{equation}
\Delta=\sum \pm a_1^{(1)}a_2^{(2)}\cdots a_n^{(n)},
\tag{8}
\end{equation}
und \gesperrt{Kronecker} noch kürzer:
\begin{equation}
\Delta=\left|a_i^{(k)}\right|.
\tag{9}
\end{equation}
Beide Bezeichnungen sind aber nur dann ganz deutlich, wenn die Elemente in der hier vorausgesetzten Weise durch zwei Indices bezeichnet sind, und durchaus unanwendbar, wenn die Elemente z.~B. numerisch gegeben sind.

Es kommen bisweilen Determinanten vor, bei denen
\[
 a_i^{(k)}=a_k^{(i)}
\]
ist, bei denen also in (1) die symmetrisch zur Diagonale des Quadrats stehenden Elemente einander gleich sind. Wir werden in diesen Fällen gewöhnlich beide Indices (um ihre Gleichwerthigkeit anzudeuten) unten hinsetzen, also
\[
 a_{i,k}=a_{k,i}
\]
setzen. Solche Determinanten heissen \gesperrt{symmetrisch}.
